\begin{table}[t]
\centering
\caption{Standardized Effect Sizes}
\label{tab:sde}
\begin{adjustbox}{max width=\textwidth}
\begin{tabular}{lccccccc}
\hline\hline
Outcome & $\hat{\beta}$ & SE & SD($X$) & SD($Y$) & SDE & SE(SDE) & Classification \\
\hline
\multicolumn{8}{l}{\textit{Panel A: Pooled}} \\
Log employment & -1.0909 & 0.421 & 0.0454 & 1.0757 & -0.0461 & 0.0178 & Small negative \\
\hline
\multicolumn{8}{l}{\textit{Panel B: Heterogeneous (by sample split)}} \\
Log emp. (private sector) & -1.5373 & 0.6463 & 0.0455 & 1.0362 & -0.0675 & 0.0284 & Moderate negative \\
Log emp. (excl.\ hospitality) & -0.9754 & 0.4244 & 0.0454 & 1.1059 & -0.04 & 0.0174 & Small negative \\
\hline\hline
\end{tabular}
\end{adjustbox}
\begin{tablenotes}[flushleft]
\small
\item \textit{Notes:} \textbf{Country:} Belgium. \textbf{Research question:} Does automatic wage indexation --- which forced 9--12\% cumulative mandatory wage increases during the 2022--2023 energy crisis --- reduce sectoral employment? \textbf{Policy mechanism:} Belgium's health-index system triggers automatic 2\% wage increases for all workers when the smoothed consumer price index crosses a preset threshold; five crossings occurred between February 2022 and September 2023, but pre-committed joint committee agreements caused different sectors to absorb these increases on different schedules (immediately, quarterly, or annually in January). \textbf{Outcome definition:} Log quarterly employment in thousands of persons by NACE Rev.~2 section, measuring total headcount of employed persons aged 15--74. \textbf{Treatment:} Continuous --- cumulative mandatory wage increase (in log points) applied to each sector through its joint committee indexation regime by each quarter. \textbf{Data:} Eurostat Labour Force Survey (\texttt{lfsq\_egan2}), Belgium, 2018-Q1 to 2024-Q4, sector--quarter panel. \textbf{Method:} Two-way fixed effects (sector and quarter FE), standard errors clustered by NACE section. \textbf{Sample:} 19 NACE sections (excluding households and extraterritorial organizations), approximately 16 pre-treatment quarters (2018-Q1 to 2021-Q4). SDE $= \hat{\beta} \cdot \text{SD}(X) / \text{SD}(Y)$ where SD($X$) is the standard deviation of the continuous treatment variable and SD($Y$) is the pre-treatment standard deviation of log employment across sectors. Classification refers to magnitude, not statistical significance: Large ($|$SDE$| > 0.15$), Moderate ($0.05$--$0.15$), Small ($0.005$--$0.05$), Null ($< 0.005$).
\end{tablenotes}
\end{table}
